% 11/18 abstract, updated with the DPO and candidate-generator results (11/20) and the information-gain / robustness runs (11/23).
\begin{abstract}
사람마다 로봇이 군중 속을 어떻게 지나가기를 원하는지는 다르다. 어떤 사람은 로봇이 일찍 양보하기를, 다른 사람은 보조를 맞추기를 선호하며,
어떤 사람은 로봇이 자신의 뒤로, 다른 사람은 앞으로 지나가기를 원한다. 기존 군중 내비게이션 정책은 고정된 사회적 준수 개념을
최적화하므로 이러한 개인별 선호에 적응하지 못한다. 우리는 소수의 쌍대(A/B) 질문(query)으로부터 사용자의 보행 스타일 선호를
식별하고, 이를 테스트 시점에 로봇의 경로 선택에 활용하는 프레임워크를 제시한다. 사람 보행자 궤적으로 학습한 VQ-VAE는
보행 행동의 코드북을 제공하며, 이로부터 스타일이 다양하고 사람다운 후보 경로를 생성한다.
안전하지 않은 후보는 걸러내고, 남은 각 후보는 해석 가능한 스타일 특징, 즉 지나가는 쪽, 앞 또는 뒤로 지나가기, 양보,
우회 폭, 사람다움으로 기술된다. 이 특징들에 대한 베이지안 Bradley--Terry 모델을 매 답변 후 갱신하며,
질문은 최선 선택지의 기대 효용(EUBO)으로 선택한다. 시뮬레이션 사용자를 이용한 Grand Central과 ETH/UCY의 학습에 쓰지 않은(held-out) 상황에서,
다섯 번의 질문은 정규화된 단순 regret을 0.45--0.47(인간 데이터 사전분포)에서 0.09--0.12로 줄이며,
네 가지 설정 모두에서 무작위 질문과 같은 예산의 DPO 학습 추정기보다 우수하고 정보 이득 질문과는 대등하다.
무작위 질문 대비 이점은 시뮬레이션 사용자가 노이즈 섞인 답을 하거나 비선형 선호를 가질 때에도 유지된다. 공통의
Grand Central 장면에서, 다섯 번의 질문 후 코드북 후보는 0.11의 regret을 남기는 반면 flow-matching 군중 내비게이션
생성기의 후보는 0.35를 남기며, 후자의 더 좁은 후보 집합은 어떤 선호로도 도달할 수 있는 수준을 제한한다. 폐루프 리플레이에서
학습된 선호를 따르면 충돌률을 높이지 않으면서 사용자의 스텝 regret을 낮게 유지한다.
\end{abstract}
